\begin{longtable}{
    %|>{\raggedright\arraybackslash}p{0.27\textwidth}
    |>{\raggedright\arraybackslash}p{0.23\textwidth}
    |p{0.55\textwidth}
    |p{0.32\textwidth}
    |p{0.27\textwidth}|
    } % 5 - 1.35 \ 
\caption{Термины цикла}
\label{tab:cycle-terms} \\
\LongTableHeader{
    \hline
    \rowcolor{black}
    %\textcolor{white}{\textbf{Русский термин}} &
    \textcolor{white}{\textbf{Термин}} &
    \textcolor{white}{\textbf{Описание}} &
    \textcolor{white}{\textbf{Источник}} &
    \textcolor{white}{\textbf{Цитата из источника}} \\
    \hline
}

\longcaptioneachpage{Источник: составлено автором}

% Цикл
\begin{minipage}[t]{\linewidth}
    \raggedright
    \begin{itemize}
        \item \textbf{цикл (рус.)}
        \item cycle (англ., фр.)
        \item zyklus (нем.)
        \item ciclo (ит., исп.)
    \end{itemize}
\end{minipage}
&
Этимология: от лат. cyclus ← греч. \textgreek{κύκλος} (круг, колесо). В экономике термин охватывает три взаимодополняющих смысла, выделенных Besomi: 1. \textquotedbl{}период времени, в течение которого характерное событие повторяется\textquotedbl{}; 2. \textquotedbl{}одно полное исполнение периодически повторяемого явления\textquotedbl{} с акцентом на порядке фаз внутри одного цикла; 3. \textquotedbl{}повторяющийся круговорот событий, регулярный порядок, серия, возвращающаяся к себе\textquotedbl{}. В экономику слово вошло через Петти (1662) в связи с сельскохозяйственными урожаями; широкую популярность получило благодаря Лорду Оверстону (1837), описавшему знаменитый десятифазный цикл торговли. Обозначение \textquotedbl{}business cycle\textquotedbl{} впервые зафиксировано у Джонсона (1887), а устойчивым стандартным термином стало у Митчелла (1913). Немецкий эквивалент konjunktur (введён Лассалем) акцентирует независимость рыночных условий от воли агентов. Итальянский ciclo economico, французский cycle économique полностью параллельны английскому.
&
\begin{minipage}[t]{\linewidth}
    \vspace{-7pt}
    \RaggedRight
    \begin{enumerate}
        \item Petty W. (1662). \textit{A Treatise of Taxes and Contributions}\cite{Petty-1662} (первое употребление в экономическом контексте)
        \item Mitchell W. C. (1927). \textit{Business Cycles: The Problem and Its Setting}\cite{Mitchell-1927} (стандартизация термина)
    \end{enumerate}
\end{minipage}
&
\textquotedbl{}The ordinary rent of land in corn … the average duration of the Cycle, within which Dearths and Plenties make their revolution\textquotedbl{}
\begin{flushright}
    Petty W., 1662\cite{Petty-1662}
\end{flushright}
\\ \hline

% Волна
\begin{minipage}[t]{\linewidth}
    \raggedright
    \begin{itemize}
        \item \textbf{волна (рус.)}
        \item wave (англ.)
        \item vague (фр.)
        \item welle (нем.)
        \item onda (ит.)
        \item springvloed (нид.)
    \end{itemize}
\end{minipage}
&
Этимология: от ст. англ. \textquotedbl{}wæg\textquotedbl{} (движение воды). В экономике метафора волны описывает более длинные и пологие движения, чем \textquotedbl{}цикл\textquotedbl{}: гидродинамический образ предполагает сначала медленный подъём, затем такой же плавный спад. Ключевой предшественник понятия – нидерландский экономист Ван Гелдерен (1913), использовавший слово \textquotedbl{}springvloed\textquotedbl{} (весенний прилив, половодье) для обозначения мощных длительных хозяйственных подъёмов; понятие связано с именем Кондратьева (1922), описавшего \textquotedbl{}большие циклы конъюнктуры\textquotedbl{} длиной 48–60 лет, которые в западной традиции через Шумпетера (1939) получили название \textquotedbl{}Kondratiev waves\textquotedbl{} или \textquotedbl{}long waves\textquotedbl{}. Клинов и Сидоров (2021) разграничивают \textquotedbl{}восходящую и нисходящую волну\textquotedbl{}. В немецкоязычной традиции используется также \textquotedbl{}konjunkturwellen\textquotedbl{} (конъюнктурные волны). Геигант (1975) обобщил употребление термина: от коротких волн Китчина до длинных волн Кондратьева. \textquotedbl{}Волна\textquotedbl{} в отличие от \textquotedbl{}цикла\textquotedbl{} акцентирует непрерывность и плавность движения, а не смену дискретных фаз.
&
\begin{minipage}[t]{\linewidth}
    \vspace{-7pt}
    \RaggedRight
    \begin{enumerate}
        \item Van Gelderen J. (1913). \textit{Springvloed: Beschouwingen over industrieele ontwikkeling en prijsbeweging}\cite{Van-Gelderen-1913} (первое употребление в экономическом контексте)
        \item Кондратьев Н. Д. (1925) \textit{Большие циклы конъюнктуры}\cite{Кондратьев-1925}
        \item Schumpeter J. (1939). \textit{Business Cycles: a Theoretical, Historical, and Statistical Analysis of the Capitalist Process}\cite{Schumpeter-1939}
        \item Geigant W. (1975). \textit{Lexikon der Volkswirtschaft}\cite{Geigant-1975}
    \end{enumerate}
\end{minipage}
&
\textquotedbl{}Длительные циклы конъюнктуры представляют собой отклонение реального уровня элементов капиталистической системы по отношению к равновесию этой самой системы\textquotedbl{}
\begin{flushright}
    Кондратьев Н. Д., 1925\cite{Кондратьев-1925}
\end{flushright}
\vspace{7pt}
\textquotedbl{}Wave-like oscillating movement of economic activity of various frequency, from Kitchin to Kondratiev cycles\textquotedbl{}
\begin{flushright}
    Geigant W., 1975\cite{Geigant-1975}
\end{flushright}
\\ \hline

% Качание / размах
\begin{minipage}[t]{\linewidth}
    \raggedright
    \begin{itemize}
        \item \textbf{размах / качание (рус.)}
        \item swing (англ.)
        \item schwingung (нем.)
    \end{itemize}
\end{minipage}
&
Этимология: от ст. англ. \textquotedbl{}swingan\textquotedbl{} (размахивать, раскачиваться); образ маятника или качелей. В экономике термин связан прежде всего с работой Кузнеца (1930), описавшего длительные движения длиной 15–25 лет, связанные со строительными и демографическими процессами; эти движения впоследствии получили название \textquotedbl{}Kuznets swings\textquotedbl{} или \textquotedbl{}Kuznets cycles\textquotedbl{}. Фриман в \textquotedbl{}The New Palgrave\textquotedbl{} (1987) подчёркивал дискуссионность самого существования длинных качаний. Клинов и Сидоров (2021) констатировали, что \textquotedbl{}ни волны Кондратьева, ни качания Кузнеца не являются регулярно повторяющимися чертами современной экономики\textquotedbl{}. Термин \textquotedbl{}swing\textquotedbl{} в отличие от \textquotedbl{}cycle\textquotedbl{} акцентирует асимметрию и нерегулярность: маятник может качаться с разной амплитудой и периодом. В словаре Ромер (2007) \textquotedbl{}swings in economic activity\textquotedbl{} используется как нейтральная альтернатива \textquotedbl{}business cycle\textquotedbl{} для описания подъёмов и спадов без импликации регулярности.
&
\begin{minipage}[t]{\linewidth}
    \vspace{-7pt}
    \RaggedRight
    \begin{enumerate}
        \item Kuznets S. (1930). \textit{Secular movements in production and prices their nature and their bearing upon cyclical fluctuations}\cite{Kuznets-1930} (оригинальная работа с описанием термина)
        \item Freeman C. (1987). \textit{Long Swings in Economic Growth}\cite{Freeman-1987} (лексикографическое закрепление термина)
        \item Romer C. D. (1993). \textit{Business Cycles}\cite{Romer-1993}
    \end{enumerate}
\end{minipage}
&
\textquotedbl{}It is highly unlikely that those who believe that long swings in economic life are a significant phenomenon will ever satisfy their statistical critics\textquotedbl{}
\begin{flushright}
    Freeman C., 1987\cite{Freeman-1987}
\end{flushright}
\vspace{7pt}
\textquotedbl{}Many modern economists prefer the term 'short-run economic fluctuations' to 'business cycle' … for describing the swings in economic activity\textquotedbl{}
\begin{flushright}
    Romer C. D., 1993\cite{Romer-1993}
\end{flushright}
\\ \hline

% Ритм
\begin{minipage}[t]{\linewidth}
    \raggedright
    \begin{itemize}
        \item \textbf{ритм (рус.)}
        \item rhythm (англ.)
        \item rythme (фр.)
        \item rhythmus (нем.)
        \item ritmo (ит.)
    \end{itemize}
\end{minipage}
&
Этимология: от греч. \textquotedbl{}\textgreek{ῥυθμός}\textquotedbl{} (мерное течение, такт); изначально музыкальная метафора. В экономике \textquotedbl{}ритм\textquotedbl{} используется для описания регулярного чередования фаз хозяйственного цикла – с акцентом на повторяемости рисунка, а не на геометрии замкнутого движения. Первое словарное закрепление применительно к \textquotedbl{}rhythmical changes of the economy\textquotedbl{} фиксируется у Вайнбергера в Staatslexikon (1929): автор противопоставляет новый \textquotedbl{}ритмический» подход к кризисам (теория цикла) старому подходу, трактовавшему кризисы как особые нарушения. Талхайм (1934) формулирует: \textquotedbl{}In the capitalistic era favourable and unfavourable times have taken a rhythmical form, so that one can speak of wave-like behaviour of economic life\textquotedbl{}. Хортон (1948) определяет деловые циклы как \textquotedbl{}rhythmic fluctuations in general business activities\textquotedbl{}. Термин \textquotedbl{}ритм\textquotedbl{} применяется и в более ранних работах неявно: Жуглар описывал \textquotedbl{}périodicité\textquotedbl{} кризисов, Шумпетер – \textquotedbl{}regularity\textquotedbl{} чередования фаз. У Туган-Барановского и Кондратьева в русской традиции \textquotedbl{}ритм\textquotedbl{} присутствует как описание правильного чередования подъёмов и спадов. Образ ритма подчёркивает не только повторяемость, но и тактность – равномерность промежутков между событиями.
&
\begin{minipage}[t]{\linewidth}
    \vspace{-7pt}
    \RaggedRight
    \begin{enumerate}
        \item Weinberger O. (1929). \textit{Konjunktur und Krisen}\cite{Weinberger-1929} (первое употребление)
        \item Horton D. C. (1948). \textit{Business Cycles}\cite{Horton-1948} (косвенное определение)
    \end{enumerate}
\end{minipage}
&
\textquotedbl{}The new approach interprets crises in the light of the rhythmical changes of the economy, so that the theory of crises has evolved into a theory of the cycle\textquotedbl{}
\begin{flushright}
    Weinberger O., 1929\cite{Weinberger-1929}
\end{flushright}
\vspace{7pt}
\textquotedbl{}Business cycles are rhythmic fluctuations in general business activities, characterized by a tendency of prices, profits, production wages and employment to move more or less together\textquotedbl{}
\begin{flushright}
    Horton D. C., 1948\cite{Horton-1948}
\end{flushright}
\\ \hline

% Спираль
\begin{minipage}[t]{\linewidth}
    \raggedright
    \begin{itemize}
        \item \textbf{спираль (рус.)}
        \item spiral (англ. фр., нем., ит.)
    \end{itemize}
\end{minipage}
&
Этимология: от лат. \textquotedbl{}spira\textquotedbl{} ← греч. \textquotedbl{}\textgreek{σπεῖρα}\textquotedbl{} (изгиб, завиток). В отличие от \textquotedbl{}цикла\textquotedbl{} (замкнутое движение), \textquotedbl{}волны\textquotedbl{} (квазипериодическое) и \textquotedbl{}ритма\textquotedbl{} (регулярное чередование), \textquotedbl{}спираль\textquotedbl{} описывает незамкнутое движение с накоплением: либо расходящееся (усиливающийся кризис – нисходящая долговая или дефляционная спираль), либо сходящееся (выход из кризиса через самоограничивающийся процесс). Концептуальную основу \textquotedbl{}debt-deflation spiral\textquotedbl{} заложил Фишер (1933), описавший самоподпитывающийся механизм дефляции долгов как \textquotedbl{}vicious spiral\textquotedbl{} нисходящего движения. Кейнс (1936) использовал образ \textquotedbl{}inflationary spiral\textquotedbl{} для самоподпитывающегося роста цен и зарплат. Классический послевоенный термин \textquotedbl{}wage-price spiral\textquotedbl{} – взаимное самоподпитывание роста зарплат и цен, окончательно закрепился в работах Гэлбрейта (1950), Самуэльсона и Солоу (1960). Мински (1986) развил метафору в понятии \textquotedbl{}debt deflation spiral\textquotedbl{} как механизма саморасширяющегося финансового кризиса. \textquotedbl{}Спираль\textquotedbl{} принципиально отличается от остальных четырёх базовых терминов: это единственная метафора, допускающая расходящееся движение без возвращения к исходной точке
&
\begin{minipage}[t]{\linewidth}
    \vspace{-7pt}
    \RaggedRight
    \begin{enumerate}
        \item Fisher I. (1933). \textit{The Debt-Deflation Theory of Great Depressions}\cite{Fisher-1933} (концептуальная основа термина)
        \item Chen B., Flaschel P. (2006). \textit{Keynesian Dynamics and the Wage–Price Spiral: Identifying Downward Rigidities}\cite{Chen-2006} (употребление термина)
    \end{enumerate}
\end{minipage}
&
\textquotedbl{}Unless some counteracting cause comes along to prevent the fall in the price level, such a depression as that of 1929–33 (starting with about 10 debtors to one creditor) tends to continue, going deeper, in a vicious spiral for many years\textquotedbl{}
\begin{flushright}
    Fisher I., 1933\cite{Fisher-1933}
\end{flushright}
\\ \hline

\end{longtable}